% Proved after four complete actual-file/source audits and exact precision refresh. Actual unit families only; no coflat object-atlas construction.
\subsection{Intrinsic integral unit jets and canonical-family openness}

\begin{proposition}[Proved: intrinsic finite jets of integral unit maps]
\label{D:intrinsic-integral-unit-Hom-jets}
Let $A$ be a complete Noetherian local $k$-algebra, with $k$
algebraically closed of characteristic $p>0$. Let $Y$ be
projective over $A$, and let $C$ be an effective Cartier
divisor whose support is the entire closed fibre. Assume
that every $jC$ is a proper $k$-scheme. Put $I=\mathcal O_Y(-C)$.
For an affine finite-type $k$-scheme $S=\operatorname{Spec}B$,
use the completion of $A\otimes_k B$ along $\mathfrak m_A$,
and the corresponding proper base change of $Y$.

Suppose that two actual length-$f$ cycles of vector bundles
$F_a,G_a$, with indices modulo $f$, are given on this
completed base change. Their arrows are meromorphic
isomorphisms
\[
 \phi_a:F_Y^*F_a\dashrightarrow F_{a+1},\qquad
 \psi_a:F_Y^*G_a\dashrightarrow G_{a+1},
                                                        \tag{D.IU.1}
\]
regular away from $C$, where $F_Y$ is absolute $p$-Frobenius,
including coefficient Frobenius. Assume a fixed bound $bC$
for every arrow and its inverse. Set
$H_a=\mathcal Hom(F_a,G_a)$ and choose $c\ge0$ such that
\[
 T_a(\eta)=\psi_a\,F_Y^*(\eta)\,\phi_a^{-1}
       \quad\hbox{maps }H_a\hbox{ into }H_{a+1}(cC).
                                                        \tag{D.IU.2}
\]
Take $c$ to work also for the reverse comparisons and the
two endomorphism cycles; $c=2b$ always does. Choose an integer $N\ge1$ with
\[
                         pN-c>N.                       \tag{D.IU.3}
\]
Then the functor of integral formal unit morphisms between
these cycles is represented by an affine scheme of finite
presentation over $S$, using only their restrictions to a
fixed finite thickening. More precisely, an integral unit
morphism is uniquely determined by the tuple
$\eta_a\in H^0(NC,H_a|_{NC})$. The exact finite equations are
\[
 \iota_N(\eta_{a+1})=T_{a,N}(\eta_a)
 \quad\hbox{in }H^0((N+c)C,H_{a+1}(cC)|_{(N+c)C}).
                                                        \tag{D.IU.4}
\]
Here $\iota_N$ is induced by $H_{a+1}\hookrightarrow
H_{a+1}(cC)$, and $T_{a,N}$ is obtained by Frobenius
pullback to $pNC$, composition with the actual pole maps,
and restriction to $(N+c)C$.

The functor of integral formal unit isomorphisms is also
affine of finite presentation: retain a reverse tuple with
the reverse unit equations and require the two products
to be the identity on $NC$. These descriptions are functorial
over every coefficient $B$-algebra, in the category of formal
bundles. Over finite-type coefficient algebras and finite
local Artinian coefficient algebras, the formal maps
algebraize on the proper Noetherian completed bases by
Grothendieck existence. No arbitrary-ring algebraization
or global lifting of an arbitrarily prescribed ordinary
Hom jet is asserted.
\end{proposition}

\begin{proof}
Every coefficient base change in the first part is interpreted
on the formal completion along $C$. Thus it is defined by the
compatible restrictions to the proper finite-level thickenings $jC$,
also when the coefficient algebra is not Noetherian. The
Cartier equation remains a nonzerodivisor: the finite
thickenings and their layer sequences are obtained by flat
coefficient extension over the field $k$.

We first type (D.IU.4) without dividing on a nonreduced
scheme. For a Cartier ideal,
\[
                 I^{N+c}H_{a+1}(cC)=I^N H_{a+1}.
                                                        \tag{D.IU.5}
\]
Consequently the inclusion of lattices induces a canonical
injection
\[
 H_{a+1}/I^N H_{a+1}\longrightarrow
 H_{a+1}(cC)/I^{N+c}H_{a+1}(cC).
                                                        \tag{D.IU.6}
\]
This is $\iota_N$. Frobenius sends $I^N$ into $I^{pN}$,
so an $N$-jet of a map determines its Frobenius pullback
on $pNC$. The pole maps in (D.IU.2), interpreted as honest
maps into the indicated twists, give a section of
$H_{a+1}(cC)$ on that thickening. Condition (D.IU.3)
implies $pN\ge N+c$, so restriction defines $T_{a,N}$.
There is therefore no arbitrary choice of a residual
fraction in (D.IU.4).

\emph{Local correction and gluing.}
Work on a common formal affine trivialization of the finitely
many bundles, and choose any integral formal lifts
$\eta_a^{(0)}$ of a solution of (D.IU.4). Such local lifts
exist by lifting matrix entries. Regard $T=(T_a)_a$ as the
cyclic operator on the tuple, so its output at phase $a+1$
uses the input at phase $a$. By (D.IU.5), the finite equation
says exactly
\[
                T\eta^{(0)}-\eta^{(0)}\in I^N\prod_a H_a.
                                                        \tag{D.IU.7}
\]
In particular the first iterate is integral. The operator
is additive and sends a difference of order at least $j$
to one of order at least $pj-c$. Thus for
$\eta^{(v+1)}=T\eta^{(v)}$ the successive differences
have order at least
\[
       N_v=p^vN-c\frac{p^v-1}{p-1}\longrightarrow\infty.
                                                        \tag{D.IU.8}
\]
Every iterate is integral and preserves the prescribed
$N$-jet. Completeness gives an integral limit $\eta$ with
$T\eta=\eta$. If two integral fixed tuples have the same
$N$-jet, their difference lies in every $I^{N_v}H_a$ and
is zero by separatedness. On an overlap two choices of
local lifts therefore give the same fixed tuple. The local
limits glue to the unique global formal unit morphism.
The construction and this uniqueness commute with every
coefficient change. It uses local formal lifts, not a
surjectivity assertion for global ordinary Hom sections.

For reverse tuples the same argument applies. The product
of two exact unit morphisms is an exact unit endomorphism.
If its $N$-jet is the identity, the just-proved uniqueness
makes it the identity formally. Thus the two finite inverse
equations are equivalent to an actual unit isomorphism.

\emph{Finite presentation.}
For each phase, the $N$-jet Hom functor is affine of finite
presentation over $S$: the finite cycle is proper and flat
over $S$, and both restricted bundles are finitely presented
and flat over $S$. The same statement applies to the target
Hom functor on $(N+c)C$, with its fixed line twist. These
are precisely the hypotheses of
\citet[Proposition 99.3.10, Tag 08K6]{StacksIntrinsicUnitHomD}.
Absolute Frobenius, the fixed pole maps and restriction
define natural polynomial transformations of these finite
Hom functors. Hence (D.IU.4) is an equality of two morphisms
between affine finite-presentation schemes. Its zero scheme
is again affine of finite presentation. Composition and the
two inverse equations are morphisms of the same finite
proper-cycle Hom functors, giving the asserted Isom scheme,
as in the paired-inverse construction of
\citet[Proposition 99.4.3, Tag 08K9]{StacksIntrinsicUnitHomD}.
This represents the actual integral formal functor by the
local contraction, rather than the image of a scheme of
uncontrolled higher-jet witnesses.

Finally, if the coefficient algebra is finite type over $k$
or finite local Artinian, the completed base ring is
Noetherian. The $C$-adic and closed-fibre topologies are
cofinal. The full faithfulness in Grothendieck existence
therefore identifies the formal maps with algebraic maps
on the proper base change of $Y$
\citep[Theorem 30.27.1, Tag 088E]{StacksIntrinsicUnitHomD}.
Their unit and inverse identities algebraize as well.
This is the only algebraization assertion used here.
\end{proof}

\begin{proposition}[Proved: specified Artinian unit lifts for canonical families]
\label{D:canonical-family-unit-Isom-openness}
In Proposition~\ref{D:intrinsic-integral-unit-Hom-jets},
suppose now that $A$ is a complete normal local surface and
$Y$ is a projective regular resolution, an isomorphism on
its puncture $U$. Assume that the actual parameter families
have the following property at every finite local Artinian
coefficient point, including every subsequent Artinian
coefficient change. For each phase there is a finite,
coefficient-flat module $M_{a,T}$ with maximal
Cohen--Macaulay closed fibre such that
\[
 M_{a,T}=H^0(U_T,F_{a,T}|_{U_T}),\qquad
 F_{a,T}=(\pi_T^*M_{a,T})^{\vee\vee},
                                                        \tag{D.IU.9}
\]
and these module identifications commute with the specified
coefficient changes. Require the same condition for every
other family being compared. The coflat, split-boundary
families in Proposition~\ref{D:weak-cycle-Artinian-full-realization}
supply (D.IU.9) whenever actual compatible unit arrows have
also been constructed; those arrows are not supplied by
that realization proposition alone.

Every prescribed unit isomorphism over a quotient of a
finite local Artinian coefficient algebra then lifts uniquely,
and its lift is integral on $Y_T$. For a fixed geometric
unit object, its integral unit-Isom scheme over $S$ is etale.
Consequently its isomorphism-class image is open and closed,
and
\[
 \#\{\hbox{unit classes represented by }S(k)\}
                         \le \#\operatorname{Irr}(S).
                                                        \tag{D.IU.10}
\]

For ordinary canonical cycles on the Cartier models of the
hereditary budget theorems, the forward arrows are regular
and the inverse arrows have pole at most $pC$. In comparing
two such cycles one may therefore take $c=p$ and $N=3$ in
(D.IU.2)--(D.IU.4), independently of the cover degree,
rank, number of phases and nilpotence order. This is an
isomorphism criterion for already actual families. It does
not construct a degree-linear parameter scheme of all such
families or remove the retained coflat/split-boundary
conditions.
\end{proposition}

\begin{proof}
Let $T\to T_0$ be a nilpotent quotient of finite local
Artinian coefficient algebras, and let a unit isomorphism
over $T_0$ be specified. On $U_T$, a length-$f$ ordinary
unit cycle is equivalent to its first bundle with a
$q=p^f$ Frobenius unit structure. The fixed-frame
construction gives its finite etale
$\operatorname{GL}_r(\mathbb F_q)$-torsor, as in the
earlier local unit and unit-Isom proofs. In local frames
this is the usual Lang equation, with invertible derivative;
the finite etale frame construction retains the maps and
their compositions. Finite etale covers and all their
morphisms are invariant under nilpotent thickenings
\citep[Lemma 58.8.3, Tag 0BQB]{StacksIntrinsicUnitHomD}.
Thus the specified unit isomorphism on $U_{T_0}$ lifts
uniquely to $U_T$.

The first identity in (D.IU.9) extends this punctured map
to an $A_T$-linear isomorphism $M_{a,T}\to M'_{a,T}$ at
every phase. Apply pullback and double dual to obtain an
integral isomorphism $F_{a,T}\to G_{a,T}$ on $Y_T$.
The base-change clause in (D.IU.9) makes its reduction
the prescribed isomorphism. It commutes with the unit
arrows: the two meromorphic maps agree on $U_T$, and the
fixed Cartier equation is a nonzerodivisor on their
locally free targets after a pole twist. Hence they agree
everywhere. The same observation proves uniqueness of
the integral extension. This is where canonical relative
fullness enters. Bounded poles of the unit arrows alone
would not provide this integral nilpotent lift.

The first proposition represents these integral unit
isomorphisms by a finite-presentation scheme. We have just
proved the required unique lifts for every finite local
Artinian extension at each closed point, not merely for
dual numbers. The local Artinian smoothness criterion,
together with uniqueness and vanishing of relative
tangent spaces, makes the morphism etale at those points;
the open etale locus is the whole scheme since it is
finite type over the algebraically closed field. These
are the exact criteria in
\citep[Tags 02HT, 02HM and 00UR]{StacksIntrinsicUnitHomD}.

Fixing a $k$-valued unit object, its Isom scheme has open
image in $S$. Such images for distinct classes are
disjoint: a nonempty intersection is open in a finite-type
$k$-scheme and has a closed $k$-point, which would identify
the two fixed objects. The union of these images contains
every closed point, hence all of $S$. Each complement is
the union of the other open images, so each class image
is also closed. Every nonempty class image contains an
irreducible component, giving (D.IU.10).

For the last assertion use the actual canonical Frobenius
sandwich of the Cartier-generation theorem. Its relative
pole statement can also be checked directly. The forward
punctured map extends to the canonical target module by
relative Hartogs, and then to the full bundles by generic
evaluation and codimension-two Hartogs. At each Artinian
coefficient base this uses the stated coefficient-flat maximal
Cohen--Macaulay target modules and their relative depth-two
Hartogs property; no regularity of a Frobenius source module
is presumed. The source bundle
$F_{a,T}$ has factorizations of $x\operatorname{Id}$ and
$y\operatorname{Id}$ through free bundles. Frobenius gives
such factorizations for $x^p$ and $y^p$ on $F_Y^*F_{a,T}$.
Compose the inverse punctured arrow with the first free-target
map. Relative module Hartogs and full pullback extend that
composite from the canonical target bundle. The second map
of the factorization then shows that both $x^p\phi_a^{-1}$
and $y^p\phi_a^{-1}$ are regular. Since $x^p,y^p$ generate
$\mathcal O_Y(-pC)$, the inverse has pole at most $pC$.
The same argument works at every Artinian coefficient base.
For a finite-type parameter base, these bounds hold on the
actual family as well: relative to its initially fixed finite
pole bound, each obstruction to the sharper bound is a
coherent section on a fixed proper finite-level thickening.
It vanishes after every Artinian coefficient specialization;
Krull intersection at the closed coefficient points therefore
makes it zero. Thus the target
forward arrow in (D.IU.2) has no pole and the source inverse
has pole at most $pC$, so $c=p$. For every prime $p\ge2$,
$3p-p>3$. Thus $N=3$ satisfies (D.IU.3), and the finite
equations are on $(p+3)C$. The integral Artinian comparison
just proved prevents any increase of this depth with the
nilpotence order.
\end{proof}

\begin{remark}[The unresolved quantitative step]
The finite Hom schemes here provide qualitative finite
presentation and exact specified-isomorphism lifting for
actual families. No coordinate or degree estimate for
cohomology base-change strata or split-boundary conditions
is inferred. The known $O(ftr^2)$ finite-cycle quotient/Hom
budgets and the present openness theorem therefore do not,
by themselves, give an $O(ftr^2)$ effective object atlas or
a hereditary representation bound.
\end{remark}
